\documentclass[11pt]{article}

\usepackage[a4paper,margin=2.5cm]{geometry}
\usepackage{amsmath}
\usepackage{amssymb}
\usepackage{mathtools}
\usepackage{booktabs}
\usepackage{xcolor}

\newcommand{\R}{\mathbb{R}}
\newcommand{\Gr}{\mathcal{G}}
\newcommand{\Ti}{\mathcal{T}}

\title{Case 1: an FCNN surrogate for the SPE11B pairwise distance}
\author{}
\date{}

\begin{document}
\maketitle

\section{Problem}

The SPE11B benchmark compares submissions through a pairwise distance table:
for every reporting year, the Wasserstein-1 distance between the CO$_2$ mass
images of each two submissions. We train a network
%
\begin{equation}
  f_{\boldsymbol\theta} : \R^{2 \times 120 \times 840} \to \R,
  \qquad
  f_{\boldsymbol\theta}(M, M') \approx d_t(M, M'),
  \label{eq:problem}
\end{equation}
%
to reproduce that table from the two images alone, as a cheap surrogate for
evaluating $d_t$.

\section{Notation}

\begin{center}
\begin{tabular}{ll}
  \toprule
  $\Gr = \{g_0 < \dots < g_{33}\}$ & the $G = 34$ submissions, alphabetically \\
  $\Ti = \{0, 5, \dots, 1000\}$ & the $T = 201$ reporting years \\
  $M \in \R^{120 \times 840}$ & one image: a field of CO$_2$ mass, in kg per cell \\
  $D^{(t)} \in \R^{34 \times 34}$ & ground-truth distance table at year $t$ \\
  \bottomrule
\end{tabular}
\end{center}

\noindent
$D^{(t)}$ is symmetric with zero diagonal, and exists for $t \neq 0$ only:
$200$ of the $201$ years.\footnote{One table per file
\texttt{dense/spe11b\_co2mass\_w1\_diff\_$t$y.csv}. Their row and column labels
drop the trailing \texttt{1} of single-submission groups
(\texttt{calgary} for $\texttt{calgary1}$); \texttt{\_distance\_lookup} undoes
the relabelling.}

\section{Data}

\paragraph{Images.} A single archive\footnote{%
\texttt{spe11b\_tmco2\_dt50y\_images.npz}, built by
\texttt{process\_map\_files.py}. The suffix \texttt{dt50y} is misleading: the
interval in the file is $5$ years.} holds one image per column, flattened
row-major over the $840 \times 120$ grid, its column index $k$ being
\texttt{index\_image\_in\_archive} in the code:
%
\begin{equation}
  A \in \R^{100800 \times 6833},
  \qquad
  M_k \coloneqq A_{:,k} \text{ reshaped},
  \qquad
  (M_k)_{r,c} = A_{\,840 r + c,\; k}.
  \label{eq:images}
\end{equation}
%
Rows run bottom-up ($r = 0$ is the deepest layer), so plotting requires
\texttt{origin="lower"}.

\paragraph{Column labels.} The companion file
\texttt{spe11b\_tmco2\_dt50y\_indices.json} is a list of labels, not a numeric
array:
%
\begin{equation}
  \mu : \{0, \dots, 6832\} \to \Gr \times \Ti,
  \qquad
  \mu(k) = (g_k, \tau_k).
  \label{eq:labels}
\end{equation}
%
It is a bijection onto $(\Gr \times \Ti) \setminus \{(\texttt{kfupm1}, 815)\}$,
the one image missing from the source data; hence the column count
%
\begin{equation}
  6833 = G\,T - 1 = 34 \cdot 201 - 1 .
\end{equation}

\paragraph{Ordering.} Columns are sorted by year, then alphabetically by group.
With $\delta_k \coloneqq \mathbf{1}[\,k \ge 5551\,]$ correcting for the missing
image,
%
\begin{equation}
  \tau_k = 5 \bigl\lfloor (k + \delta_k) / G \bigr\rfloor,
  \qquad
  g_k = g_{\,(k + \delta_k) \bmod G} .
  \label{eq:ordering}
\end{equation}
%
The year blocks $\mathcal{I}_t \coloneqq \{k : \tau_k = t\}$ are therefore
contiguous, of size $34$ ($33$ for $t = 815$):
%
\begin{equation}
  \mathcal{I}_{5s} = \{34 s, \dots, 34 s + 33\},
  \qquad s \le 162 .
  \label{eq:blocks}
\end{equation}

\section{Dataset}

The four steps below build the training pairs. Their symbols carry the
following names in \texttt{utils\_datasets.py}:
%
\begin{center}
\begin{tabular}{ll}
  \toprule
  $k$ & \texttt{index\_image\_in\_archive} \\
  $\pi(k)$ & \texttt{index\_image} \\
  $\ell$, a row of $\mathcal{P}$ & \texttt{index\_pair} \\
  $\mathcal{P}$ & \texttt{index\_image\_pairs} \\
  \bottomrule
\end{tabular}
\end{center}

\paragraph{1. Column selection.} \texttt{create\_datasets} keeps
$\mathcal{K} \coloneqq \texttt{range(start, stop, step)}$, $n \coloneqq
|\mathcal{K}|$, and re-indexes it by position,
%
\begin{equation}
  \pi(k) = \frac{k - \texttt{start}}{\texttt{step}} \in \{0, \dots, n-1\},
  \qquad
  X \in \R^{\,n \times 120 \times 840},
  \quad X_{\pi(k)} = M_k .
  \label{eq:selection}
\end{equation}
%
In \texttt{main.py}, $\texttt{start} = 34 = G$ drops exactly the year-$0$ block
$\mathcal{I}_0$, which has no $D^{(t)}$, and $\texttt{step} = 1$, so
$\pi(k) = k - 34$ and $n = 6799$.

\paragraph{2. Pairing, within a year.} With $\mathcal{K}_t \coloneqq
\mathcal{K} \cap \mathcal{I}_t$ and $n_t \coloneqq |\mathcal{K}_t| \le G$,
%
\begin{equation}
  \mathcal{P} \coloneqq
  \biguplus_{t} \; \pi(\mathcal{K}_t) \times \pi(\mathcal{K}_t),
  \qquad
  N \coloneqq |\mathcal{P}| = \sum_t n_t^2 .
  \label{eq:pairs}
\end{equation}
%
Both orderings and the diagonal are included; no pair crosses years.

\begingroup\color{red}
\noindent\textbf{Why within a year: a constraint of the data, not a modelling
choice.} The benchmark publishes one $34 \times 34$ table per year, so
$D^{(t)}_{pq}$ is defined only for two submissions compared \emph{at the same}
$t$. No ground-truth distance exists between images at different times, e.g.\
between $(\texttt{calgary1}, 5\,\mathrm{y})$ and
$(\texttt{sintef1}, 900\,\mathrm{y})$. Pairing across timesteps would require
computing the Wasserstein-1 distance ourselves, which is done nowhere in this
repository.
\endgroup

\paragraph{3. Targets.} For $\ell = (\pi(i), \pi(j)) \in \mathcal{P}$,
%
\begin{equation}
  y_\ell \coloneqq D^{(\tau_i)}_{\,p(i),\,p(j)},
  \qquad
  p(k) \coloneqq (k + \delta_k) \bmod G
  \ \ \text{the rank of } g_k \text{ in } \Gr .
  \label{eq:targets}
\end{equation}

\paragraph{4. Scaling and splitting.} Inputs and targets are min--max scaled to
$[0,1]$ independently, with extrema taken over the whole dataset; the rows of
$\mathcal{P}$ are then cut into contiguous $70/15/15$ slices, unshuffled.

\paragraph{Index bounds.} A row of $\mathcal{P}$ has both entries in
$\{0, \dots, n-1\}$, but within one year block they range only over the $n_t$
consecutive positions of that block. For \texttt{main.py} the block of year
$5s$ occupies rows $1156(s-1), \dots, 1156s - 1$ with entries in
$\{34(s-1), \dots, 34(s-1)+33\}$, so the array opens with
%
\begin{equation}
  (0,0), (0,1), \dots, (0,33),\ (1,0), \dots, (33,33),\ (34,34), \dots
\end{equation}
%
The bound $33$ is $G - 1$: the submissions available at one year, not a number
of timesteps. A single submission contributes $T = 201$ images, but spread over
$201$ different blocks.

\paragraph{Counts.} For \texttt{main.py}, $n_t = 34$ for all years but
$n_{815} = 33$, so
%
\begin{equation}
  N = 199 \cdot 34^2 + 33^2 = 231133
  \;\longrightarrow\;
  161793 \,/\, 34669 \,/\, 34671 .
  \label{eq:counts}
\end{equation}
%
Two consequences of the choices above are worth stating. Scaling before
splitting makes the two extrema statistics of the test set as well. And since
the rows of $\mathcal{P}$ are ordered by year, a contiguous split is a
\emph{temporal} one: training covers years $5$--$700$, validation
$700$--$855$, test $855$--$1000$, so the model is assessed by extrapolation in
time.

The three sets share one image stack $X$; only $\mathcal{P}$ is split, so
memory grows with $n$ and not with $N$.

\section{Model}

A pair is flattened to $x_\ell \in \R^{d}$, $d = 2 \cdot 120 \cdot 840 =
201600$, and passed through a fully connected network of widths $[d, 64, 64,
1]$. Activations are SiLU, $\mathrm{silu}(x) = x\,\sigma(x)$, on every layer but
the last, which is linear. Weights use He initialization,
$W \sim \mathcal{N}(0, 2/\mathrm{fan\_in})$; biases start at zero.

\section{Training}

Adam on the mean squared error over minibatches $B$,
%
\begin{equation}
  \mathcal{L}(\boldsymbol\theta)
  = \frac{1}{|B|} \sum_{\ell \in B}
    \bigl( f_{\boldsymbol\theta}(x_\ell) - y_\ell \bigr)^2 ,
  \label{eq:loss}
\end{equation}
%
with the training rows reshuffled each epoch. \texttt{main.py} uses
$|B| = 128$, learning rate $10^{-2}$ and $100$ epochs, i.e.\ $1264$ steps per
epoch; losses are logged every $10$.

\section{Test metrics}

Let $e_i \coloneqq \hat{y}_i - y_i$ be the residual on the $i$-th of the $N$
test pairs and $\bar{y}$ the mean target. \texttt{postprocess.py} reports
%
\begin{align}
  \mathrm{NMSE} &= \frac{\sum_i e_i^2}{\sum_i y_i^2},
  &
  \mathrm{NMAE} &= \frac{\sum_i |e_i|}{\sum_i y_i},
  \label{eq:nmse} \\[4pt]
  \mathrm{NRMSE} &= \sqrt{\mathrm{NMSE}},
  &
  R^2 &= 1 - \frac{\sum_i e_i^2}{\sum_i (y_i - \bar{y})^2} .
  \label{eq:r2}
\end{align}
%
The first three are each normalized by a statistic of the targets, hence
dimensionless; they are \emph{not} the plain MSE, RMSE and MAE, despite the
names \texttt{mse}, \texttt{rmse}, \texttt{mae} carrying them in the code. When
$\sum_i (y_i - \bar{y})^2 = 0$, the convention is $R^2 = 1$ if the residuals
vanish and $R^2 = 0$ otherwise.
\texttt{plot\_r\_squared} scatters $\hat{y}_i$ against $y_i$ with the identity
line and annotates $R^2$.

\end{document}
